\section{Chapter Exercise Bank}
\label{sec:chapter1-exercise-bank}

Unless stated otherwise, use a right-handed Cartesian basis and give exact
answers before decimal approximations.

\subsection*{Foundations}

\begin{chapterexercise}[Scalar or vector?]\difficulty{1}
Classify each quantity as scalar or vector and justify the classification:
energy, angular momentum, electric potential, magnetic field, probability
density, and probability current.
\end{chapterexercise}

\begin{chapterexercise}[Components and magnitude]\difficulty{1}
For $\mathbf A=2\mathbf e_x-3\mathbf e_y+6\mathbf e_z$, calculate
$\lVert\mathbf A\rVert$ and $\widehat{\mathbf A}$.
\end{chapterexercise}

\begin{chapterexercise}[Vector addition]\difficulty{1}
A hiker walks $5\,\mathrm{km}$ east and then $12\,\mathrm{km}$ north. Write the
resultant displacement in component form and find its magnitude.
\end{chapterexercise}

\begin{chapterexercise}[Dot product and angle]\difficulty{1}
For $\mathbf A=(1,2,2)$ and $\mathbf B=(2,1,-2)$, calculate
$\mathbf A\cdot\mathbf B$ and the angle between the vectors.
\end{chapterexercise}

\subsection*{Intermediate}

\begin{chapterexercise}[Projection]\difficulty{2}
Resolve $\mathbf A=(4,3,0)$ into components parallel and perpendicular to
$\mathbf B=(1,1,0)$.
\end{chapterexercise}

\begin{chapterexercise}[Torque]\difficulty{2}
A force $\mathbf F=(0,8,0)\,\mathrm{N}$ is applied at
$\mathbf r=(0.25,0,0)\,\mathrm{m}$. Calculate
$\boldsymbol\tau=\mathbf r\times\mathbf F$ and interpret its direction.
\end{chapterexercise}

\begin{chapterexercise}[Area of a triangle]\difficulty{2}
The vertices of a triangle are $P=(0,0,0)$, $Q=(2,1,0)$, and $R=(1,3,0)$.
Use a cross product to find its area.
\end{chapterexercise}

\begin{chapterexercise}[Change of basis]\difficulty{2}
Let
\[
\mathbf e'_1=\frac{\mathbf e_x+\mathbf e_y}{\sqrt2},\qquad
\mathbf e'_2=\frac{-\mathbf e_x+\mathbf e_y}{\sqrt2}.
\]
Find the components of $\mathbf v=3\mathbf e_x+\mathbf e_y$ in the primed basis
and verify that its magnitude is unchanged.
\end{chapterexercise}

\subsection*{Advanced}

\begin{chapterexercise}[Linear independence]\difficulty{3}
Determine all values of $\lambda$ for which the vectors
$(1,0,1)$, $(0,1,1)$, and $(1,1,\lambda)$ are linearly dependent.
\end{chapterexercise}

\begin{chapterexercise}[Orthonormalization]\difficulty{3}
Apply the Gram--Schmidt procedure to $\mathbf v_1=(1,1,0)$ and
$\mathbf v_2=(1,0,1)$ to construct an orthonormal basis for their span.
\end{chapterexercise}

\begin{chapterexercise}[Coordinate-free proof]\difficulty{3}
Starting from the dot-product axioms, prove the Pythagorean theorem for vectors:
if $\mathbf A\cdot\mathbf B=0$, then
$\lVert\mathbf A+\mathbf B\rVert^2
=\lVert\mathbf A\rVert^2+\lVert\mathbf B\rVert^2$.
\end{chapterexercise}

\begin{chapterexercise}[From arrows to functions]\difficulty{3}
Show that the functions $1$, $x$, and $x^2$ are linearly independent on any
interval containing more than two points. Explain why this demonstrates that
vectors need not be geometric arrows.
\end{chapterexercise}

\begin{bookinsight}[Self-check]
A complete solution should state the basis, preserve units, distinguish a
vector from its components, and interpret the final direction or sign whenever
it has physical meaning.
\end{bookinsight}

\subsection*{Additional Foundations}
\begin{chapterexercise}[Units and components]\difficulty{1}
A velocity is written as $\mathbf v=(4,-3,0)\,\mathrm{m\,s^{-1}}$. State the units of each component, the magnitude, and the corresponding unit vector.
\end{chapterexercise}
\begin{chapterexercise}[Displacement between points]\difficulty{1}
Find the displacement from $P=(-2,1,4)$ to $Q=(3,-1,7)$ and its length.
\end{chapterexercise}
\begin{chapterexercise}[Opposite vectors]\difficulty{1}
For what scalar $\lambda$ are $\mathbf a=(2,-4,6)$ and $\mathbf b=(-1,2,-3)$ related by $\mathbf a=\lambda\mathbf b$?
\end{chapterexercise}
\begin{chapterexercise}[Direction cosines]\difficulty{1}
Find the angles made by $\mathbf a=(2,2,1)$ with the positive coordinate axes.
\end{chapterexercise}
\begin{chapterexercise}[Triangle inequality]\difficulty{1}
Verify the triangle inequality for $\mathbf a=(1,2)$ and $\mathbf b=(-2,2)$. When does equality occur geometrically?
\end{chapterexercise}
\begin{chapterexercise}[Orthogonality test]\difficulty{1}
Determine $k$ so that $(2,k,-1)$ is perpendicular to $(1,3,4)$.
\end{chapterexercise}
\begin{chapterexercise}[Work and sign]\difficulty{1}
A force of magnitude $20\,\mathrm N$ acts through $5\,\mathrm m$ at angles $0^\circ$, $90^\circ$, and $150^\circ$. Compute the work in each case and interpret the sign.
\end{chapterexercise}
\begin{chapterexercise}[Standard basis products]\difficulty{1}
List all nine products $\mathbf e_i\cdot\mathbf e_j$ and all ordered cross products of distinct Cartesian basis vectors.
\end{chapterexercise}
\begin{chapterexercise}[Parallelogram diagonal]\difficulty{1}
Two sides of a parallelogram are $\mathbf a=(3,1)$ and $\mathbf b=(1,2)$. Find both diagonals.
\end{chapterexercise}
\begin{chapterexercise}[Unit direction]\difficulty{1}
Find the unit vector pointing from $A=(1,2,-1)$ to $B=(5,5,1)$.
\end{chapterexercise}
\begin{chapterexercise}[Scalar projection]\difficulty{1}
Find $\operatorname{comp}_{\mathbf b}\mathbf a$ for $\mathbf a=(3,4)$ and $\mathbf b=(1,0)$.
\end{chapterexercise}
\begin{chapterexercise}[Vector projection]\difficulty{1}
Project $\mathbf a=(2,5)$ onto $\mathbf b=(2,0)$ and explain why scaling $\mathbf b$ does not alter the projected vector.
\end{chapterexercise}

\subsection*{Additional Intermediate Problems}
\begin{chapterexercise}[Distance from a line]\difficulty{2}
For $\mathbf r=(4,3,0)$ and line direction $\mathbf a=(1,1,0)$ through the origin, use orthogonal decomposition to find the shortest distance from the point to the line.
\end{chapterexercise}
\begin{chapterexercise}[Inclined-plane decomposition]\difficulty{2}
Resolve a weight $m\mathbf g$ into components parallel and perpendicular to a plane inclined at angle $\alpha$ above the horizontal.
\end{chapterexercise}
\begin{chapterexercise}[Magnetic-force direction]\difficulty{2}
For a positive charge moving along $+x$ in a magnetic field along $+z$, determine the direction of $q\mathbf v\times\mathbf B$. Repeat for an electron.
\end{chapterexercise}
\begin{chapterexercise}[Angular momentum]\difficulty{2}
A particle has $\mathbf r=(2,0,0)\,\mathrm m$ and $\mathbf p=(0,3,4)\,\mathrm{kg\,m\,s^{-1}}$. Find $\mathbf L=\mathbf r\times\mathbf p$.
\end{chapterexercise}
\begin{chapterexercise}[Coplanarity]\difficulty{2}
Use the scalar triple product to test whether $(1,0,1)$, $(0,2,1)$, and $(2,2,3)$ are coplanar.
\end{chapterexercise}
\begin{chapterexercise}[Volume]\difficulty{2}
Find the volume of the parallelepiped spanned by $(1,1,0)$, $(0,2,1)$, and $(2,0,1)$.
\end{chapterexercise}
\begin{chapterexercise}[Basis coordinates]\difficulty{2}
Express $\mathbf v=(5,1)$ in the basis $\mathbf b_1=(1,1)$, $\mathbf b_2=(1,-1)$.
\end{chapterexercise}
\begin{chapterexercise}[Basis verification]\difficulty{2}
Show that $(1,1,0)$, $(0,1,1)$, and $(1,0,1)$ form a basis of $\mathbb R^3$.
\end{chapterexercise}
\begin{chapterexercise}[Matrix representation]\difficulty{2}
Find the matrix of the transformation $T(x,y)=(2x-y,x+3y)$ in the standard basis.
\end{chapterexercise}
\begin{chapterexercise}[Image of a basis]\difficulty{2}
A linear map satisfies $T\mathbf e_1=(1,2)$ and $T\mathbf e_2=(-1,3)$. Find $T(4,-2)$.
\end{chapterexercise}
\begin{chapterexercise}[Eigenvector check]\difficulty{2}
Determine whether $(1,1)$ and $(1,-1)$ are eigenvectors of $A=\begin{pmatrix}2&1\\1&2\end{pmatrix}$, and find the corresponding eigenvalues.
\end{chapterexercise}
\begin{chapterexercise}[Polynomial coordinates]\difficulty{2}
Write $p(x)=1+4x-2x^2$ in the basis $\{1,x,x^2\}$ and in the basis $\{1,1+x,(1+x)^2\}$.
\end{chapterexercise}

\subsection*{Additional Advanced Problems}
\begin{chapterexercise}[Cauchy--Schwarz]\difficulty{3}
Use $\lVert\mathbf a-t\mathbf b\rVert^2\ge0$ for all real $t$ to prove $|\mathbf a\cdot\mathbf b|\le\lVert\mathbf a\rVert\lVert\mathbf b\rVert$.
\end{chapterexercise}
\begin{chapterexercise}[Equality in Cauchy--Schwarz]\difficulty{3}
Determine the condition for equality in the Cauchy--Schwarz inequality and interpret it geometrically.
\end{chapterexercise}
\begin{chapterexercise}[Projection operator]\difficulty{3}
For a unit vector $\widehat{\mathbf n}$, define $P=\widehat{\mathbf n}\widehat{\mathbf n}^{\mathsf T}$. Prove that $P^2=P$ and interpret this property.
\end{chapterexercise}
\begin{chapterexercise}[Orthogonal complement]\difficulty{3}
Find a basis for all vectors in $\mathbb R^3$ perpendicular to $(1,2,3)$.
\end{chapterexercise}
\begin{chapterexercise}[Gram--Schmidt in three dimensions]\difficulty{3}
Apply Gram--Schmidt to $(1,1,0)$, $(1,0,1)$, and $(0,1,1)$.
\end{chapterexercise}
\begin{chapterexercise}[Change-of-basis matrix]\difficulty{3}
Construct the matrix taking standard coordinates to coordinates in the basis $(1,1)$, $(1,-1)$, and verify its inverse.
\end{chapterexercise}
\begin{chapterexercise}[Rotation preserves products]\difficulty{3}
Prove that if $R^{\mathsf T}R=I$, then $(R\mathbf a)\cdot(R\mathbf b)=\mathbf a\cdot\mathbf b$.
\end{chapterexercise}
\begin{chapterexercise}[Determinant and orientation]\difficulty{3}
Explain why an orthogonal matrix with determinant $-1$ preserves lengths but reverses orientation.
\end{chapterexercise}
\begin{chapterexercise}[Subspace test]\difficulty{3}
Determine which of the following are subspaces of $\mathbb R^3$: $x+y+z=0$; $x+y+z=1$; $xy=0$.
\end{chapterexercise}
\begin{chapterexercise}[Function-space inner product]\difficulty{3}
Using $\langle f,g\rangle=\int_{-1}^{1}f(x)g(x)\,dx$, determine whether $1$, $x$, and $x^2-1/3$ are mutually orthogonal.
\end{chapterexercise}
\begin{chapterexercise}[Eigenbasis]\difficulty{3}
Diagonalize $A=\begin{pmatrix}4&1\\1&4\end{pmatrix}$ using an orthonormal eigenbasis.
\end{chapterexercise}
\begin{chapterexercise}[Quantum-state preview]\difficulty{3}
Let $|\psi\rangle=a|0\rangle+b|1\rangle$ with complex coefficients. State the normalization condition and explain how it is the analogue of a unit vector.
\end{chapterexercise}

\subsection*{Challenge Problems}
\begin{chapterexercise}[Lagrange identity]\difficulty{4}
Prove $\lVert\mathbf a\times\mathbf b\rVert^2+(\mathbf a\cdot\mathbf b)^2=\lVert\mathbf a\rVert^2\lVert\mathbf b\rVert^2$.
\end{chapterexercise}
\begin{chapterexercise}[Best approximation]\difficulty{4}
Prove that the orthogonal projection of $\mathbf v$ onto a subspace $W$ is the unique vector in $W$ closest to $\mathbf v$.
\end{chapterexercise}
\begin{chapterexercise}[Dual viewpoint]\difficulty{4}
Explain how the dot product with a fixed vector defines a linear functional. In finite-dimensional Euclidean space, show how every linear functional arises this way.
\end{chapterexercise}
\begin{chapterexercise}[Spectral preview]\difficulty{4}
For a real symmetric $2\times2$ matrix, show directly that eigenvectors associated with distinct eigenvalues are orthogonal.
\end{chapterexercise}
